\cotefile{16}
\scanpage{3}\mspage{1}

\uline{Sur certains anneaux gradués}

\inserted{Soit $k$ anneau (\unsure{comm.|commut.}) fixé.}

1. Soi\deleted{en}t \deleted{$k$ un anneau,} $M^*$ un $k$-module gradué \inserted{par $\mathbb Z$} à graduation \inserted{finie}, $\unsure{E|\mathcal E}^*$ un\deleted{e} \deleted{algèbre graduée} \inserted{sous-anneau de l'anneau} d'endomorphismes \inserted{du module \unsure{ss-jacent} (\unsure{s/s graduation})} de $M^*$. On dénote par $\pi_i$ la projection canonique de $M^*$ sur $M^i$. \deleted{et} On suppose que l'on a
\[
 \pi_i\in E^* \qquad \text{pour tout }i
 \tag*{(\unsure{1.2|1.1})}
\]
Cela implique que $E^*$ est une \uline{sous}-\deleted{\uline{algèbre}} \inserted{anneau} graduée de l'algèbre graduée \deleted{\unsure{des}} $\operatorname{End}(M^*)$ [\deleted{\unsure{dont la}} graduation est définie, rappelons le, par
\[
 \operatorname{End}^{(\unsure{\nu|v})}(M)=\sum\operatorname{Hom}(M^i,M^{i+\nu})
\]
ou de façon équivalente
\[
 \operatorname{End}^{(\nu)}(M)=\sum_i\pi_{i+\nu}\operatorname{End}(M^*)\pi_i,
\]
la composante homogène de degré $\nu$ d'un endomorphisme $u$ de $M^*$ étant
\[
 \deleted{\unsure{u|\mathrm{gr}}}\operatorname{pr}_{\nu}(u)=\sum_i\pi_{i+\nu}u\pi_i .
\]

La graduation de $E^*$ est définie \inserted{intrinsèquement}, en termes de la \deleted{\unsure{gr}} donnée des $\pi_i$, par
\[
\left\{\begin{aligned}
 \text{\textcircled{\scriptsize a}}\quad &\boxed{E^{(\nu)}=\sum_i\pi_{i+\nu}E^*\pi_i}\\
 \text{\textcircled{\scriptsize b}}\quad &\operatorname{pr}^{E}_{\nu}u=\sum_i\pi_{i+\nu}u\pi_i
\end{aligned}\right.
\tag*{(\deleted{\unsure{1.2}}1.3)}
\]
et
\[
 u\in E^{(\nu)}\iff u\pi_i=\pi_{i+\nu}u\quad\text{pour tt } i
 \tag*{(1.4.)}
\]
\deleted{\illegible{2}}
\inserted{Noter les relations
\[
 \deleted{\text{Notons}}\quad
 \left|\begin{aligned}
 \pi_i^2&=\pi_i\\
 \pi_i\pi_j&=0\quad\text{si }i\ne j\\
 \sum\pi_i&=1
 \end{aligned}\right.
 \tag*{(1.1)}
\]}

2. Soit $E^*$ un \deleted{$k$-algèbre} \inserted{anneau} gradué, et soient $\pi_i$ ($i\in\mathbb Z$) des éléments de $E$, \inserted{nuls sauf un nb fini.} A quelle condition les formules (\unsure{1.3}) définissent-elles
% archive sheet number: 2
\scanpage{4}\mspage{2}

sur $\unsure{E|\mathcal E}^{(\unsure{\nu|v})}$ une structure d'\deleted{algèbre}anneau gradué ?

Il suffit que l'on ait les relations (1.1.)
\[
 \deleted{\unsure{*}}\left\{\begin{aligned}
 \pi_i^2&=\pi_i\\
 \pi_i\pi_j&=\pi_j\pi_i=0\quad\text{si }i\ne j\\
 \sum\pi_i&=1
 \end{aligned}\right.
\]
qui impliquent que les $\operatorname{pr}_{\nu}^{E}$ définis par (1.2.) sont des projections deux à deux orthogonales de somme 1, donc que les $E^{(\nu)}$ définissent bien une graduation \inserted{de $k$-module} sur $E^{(\nu)}$, \deleted{\unsure{puis}} \inserted{et} que $E^{(\nu)} E^{(\nu')}\subset E^{(\nu+\nu')}$. Noter aussi que (1.4.) que (1.4) reste valable, car \deleted{\unsure{comme on a pour tt} $u$} \inserted{on aura la décomposition}
% slanted
\margin{on peut remplacer $\mathbb Z$ par un \uline{groupe} quelconque}
\[
 \deleted{u=\sum_{ij}\pi_j u\pi_i}
\]
\deleted{\unsure{d'où la relation|de la relation}}
\[
 \deleted{\pi_{i+\nu}u\pi_i=\pi_{i+\nu}u\pi_i}\qquad E_i^j
\]
\[
 \left\{\begin{aligned}
 E&=\bigoplus_{i,j}\pi_j E\pi_i\\
 \operatorname{pr}_{i,j}(u)&=\pi_j u\pi_i
 \end{aligned}\right.\qquad\text{donc }\ u=\sum_{\alpha,\beta}\pi_\beta u\pi_\alpha
\]
et la relation $u\pi_i=\pi_{i+\nu}u$ implique, puisque $u\pi_i=\sum_{\beta}\pi_{\beta}u\pi_i$, et $\pi_{\unsure{j|i}+\nu}u=\sum_i\pi_{j+\nu}u\pi_\alpha$, que $\pi_\beta u\pi_i=0$ si $\beta\ne i+\nu$.

\margin{Regarder les modules sur \deleted{\unsure{un|des}} anneau $E$ muni \deleted{\unsure{de}} \deleted{\unsure{des}} d'éléments $\pi_i$ comme dessus.}

3. \deleted{\unsure{Nous notons|Nous pouvons}} Partons enfin d'une \deleted{algèbre} \inserted{anneau} graduée $E^*$ \deleted{sur $k$}, et posons nous la question de l'existence et \uline{l'unicité} d'une famille $(\pi_i)_{i\in\mathbb Z}$ d'élts de $E^*$ \inserted{satisfaisant (1.1.)}, définissant la graduation par (1.3.) (\deleted{\unsure{donc}} \inserted{ou} par (1.4)). On \deleted{\unsure{supp}} supposera ici que
% archive sheet number: 3
\scanpage{5}\mspage{3}

\deleted{\unsure{$E$ est à degrés pairs}, et que} l'on voudra que
\[
 \overline{\left|\ \pi_i=0\quad\text{si }i\notin[0,2n]\ \right|}\tag*{(3.0)}
\]
intervalle \inserted{(\unsure{pair|fini})} fixé. \deleted{\unsure{Suppo|Supp}} On s'occupera ici de l'\uline{unicité}. \deleted{Supposons établie l'unicité}

Notons d'abord que
\[
 \unsure{E|\mathcal E}^0\ni u\iff (u\pi_i=\pi_i u\qquad\forall i)
\]
donc $E^0$ est le commutant de la $k$-algèbre engendrée par les $\pi_i$. Donc les $\pi_i$ sont dans le centre $\unsure{Z|\mathfrak Z|\mathfrak z}(E^0)$ de $E^0$, et ce sont des projecteurs \inserted{mutuellement} orthogonaux de somme 1 de cet anneau commutatif. Supposons déterminé déjà les $\pi_\alpha$ pour $\alpha\notin[i,2n-i]$, \inserted{$i<n$,} et déterminons $\pi_i$ et $\pi_{2n-i}$ [i.e $\pi_\alpha$ pour $\alpha\notin[i+1,2n-i-1]$. \inserted{C'est clair si $i=n-1$, car on connaît tous $\pi_\alpha$ pour $\alpha\ne n$, d'où aussi $\pi_n=1-\sum_{\alpha\ne n}\pi_\alpha$. Reste le cas $i<n-1$} On connaît donc
\[
 \begin{aligned}
 \varphi_i&=\deleted{\unsure{(1-\sum\pi_\alpha)}}(1-\sum_{\alpha<i\deleted{\unsure{\alpha}}}\pi_i)=\sum_{\alpha\ge i}\pi_\alpha,\quad\text{et}\quad
 \psi_{2n-i}=(1-\sum_{\alpha>2n-i}\pi_i)=\sum_{\alpha\le 2n-i}\pi_{\unsure{\alpha'|\alpha}}
 \end{aligned}
\]
donc aussi la composée
\[
 \unsure{\text{Or}}\quad \psi_{2n-i}\deleted{\unsure{\varphi_i}}E^{2n-2i}\varphi_i
 =\sum_\alpha\psi_{2n-i}\pi_{\alpha+(2n-2i)}E\pi_\alpha\varphi_i)
\]
Or
\[
 \pi_\alpha\varphi_i=\begin{cases}0&\text{si }\alpha<i\\ \pi_\alpha&\text{si }\alpha\ge i\end{cases}
 \quad\text{et}\quad
 \psi_{2n-i}\pi_{\alpha+(2n-2i)}
 =\begin{cases}
 0&\text{si }\inserted{\alpha+(2n-2i)>2n-i}\ \text{i.e. }\alpha>i\\
 \pi_{\alpha+(2n-2i)}&\text{si }\alpha\le i
 \end{cases}
\]
donc le terme d'indice $\alpha$ de la dernière $\sum$ est nul sauf si $\alpha=i$, auquel cas il est égal à
\[
 \pi_{\unsure{\alpha|i}+(2n-2i)}E\pi_i=\pi_{2n-i}E\pi_i=E_{2n-i,i}\,.
\]
Donc la connaissance des $\pi_\alpha$ pour $\alpha\notin[i,2n-i]$ implique la connaissance de $E^{2n-i}_{i}\deleted{\unsure{,i}}$. \deleted{\unsure{Considér|considéré}} La condition qui va impliquer l'unicité \deleted{\unsure{est}} de $u_i$,
% archive sheet number: 4
\scanpage{6}\mspage{4}

$u_{2n-i}$ est alors la \deleted{\unsure{ques|sous}} conjonction des suivantes :
\[
\left[\begin{aligned}
 (3.1.)_i\quad &(u\in \unsure{Z|\mathfrak Z|\mathfrak z}\unsure{E|\mathcal E}^0,\ \deleted{\unsure{u\pi_\alpha=0\ \text{pour }\alpha<i},}\ 
 \text{\textcircled{$E_i^{2n-i}$}}u=0)\Longrightarrow \deleted{\unsure{u=0}}\inserted{\pi_i u=0},\\
 (3.2)_i\quad &(u\in ZE^0,\ \deleted{\unsure{\pi_\alpha u=0\ \text{pour }\alpha\ne 2n-i},}\ uE_i^{2n-i}=0)
 \Longrightarrow\deleted{\unsure{u=0}}\inserted{\pi_{2n-i}u=0}
\end{aligned}\right]
\]
En effet, la relation (3.1.) implique que \deleted{\unsure{$\pi_i(ZE^0)$|$\pi_i\in ZE^0$}} \inserted{l'annulateur de $\pi_i$ dans $ZE^0$} est l'idéal des $u\in ZE^0$ tels que $E_i^{2n-i}u=0$ ; or dans un anneau commutatif, un idempotent est connu quand on connaît son annulateur. De même pour $\pi_{2n-i}$, son annulateur dans $ZE^0$ est formé des $u\in ZE^0$ tels que $uE_i^{2n-i}=0$.

En même temps, les conditions $(3.1)_i$, $(3.2.)_i$ pour \unsure{tt|cet} $i$, donnent un \deleted{condition} procédé pour déterminer de proche en proche les $\pi_i$.

\medskip\hrule\medskip

4. Lorsque $E^*$ est comme dans 1., une condition suffisante pour avoir $(3.1)_i$ \inserted{resp. $(3.2)_i$} est la \deleted{\unsure{suff}} suivante

$(4.1)_i$ L'intersection des noyaux des $\unsure{v|u}:M^i\longrightarrow M^{2n-i}$, pour $v\in E_i^{2n-i}$, est réduite à 0.

$(4.2)_i$ La somme des images des $v:M^i\unsure{-|\to}M^{2n-i}$, pour $v\in E_i^{2n-i}$, est \deleted{\unsure{se}} égale à $M^{2n-i}$.

Ces deux conditions sont satisfaites si on a

$(4.3)_i$ Il existe \inserted{$v:M^i-M^{2n-i}$,} $v\in E_i^{2n-i}$, qui soit un \unsure{isom$^{\mathrm e}$|isom.} de
% archive sheet number: 5
\scanpage{7}\mspage{5}

façon équivalente, il existe un $v\in \unsure{E|\mathcal E}^{2n-2i}$, tel que le morphisme induit $v:M^i- M^{2n-i}$ soit un isomorphisme.

On peut transformer ceci de la façon suivante :
\[
 \left\{\begin{aligned}
 &\exists\ v_i\in E^{2n-2i},\ w_i\in E^{-(2n-2i)},\quad\text{tels que l'on ait}\\
 &(w_i v_i-1)\pi_i=0,\qquad(v_i w_i-1)\pi_{2n-i}=0
 \end{aligned}\right.
 \tag*{$(4.4.)_i$}
\]
Or, sans plus supposer que $E^*$ se réalise à l'aide d'un $M^*$, on voit aussitôt que $(4.4.)_i$ implique $(3.1)_i$ et $(3.2)_i$. En effet, quitte à remplacer $v,w$ par $v\pi_i=\pi_{2n-i}v$, $w\pi_{2n-i}=\pi_i w$, on \unsure{va supposer|les suppose|peut supposer} $v,w$ dans $E_i^{2n-i}$, $E_{2n-i}^i$. Si alors \deleted{\unsure{$Z$}} $u\in E^{(0)}$ est tel que $E_i^{2n-i}u=0$, donc $vu=0$ on aura $wvu\pi_i\deleted{\unsure{0}}=0$ i.e. $wv\pi_i u=0$ donc $\pi_i u=0$ (car $wv\pi_i=\pi_i$). De $=$ pour $(3.2)_i$, en utilisant la deuxième relation $(\unsure{4.4}.)_i$.

\margin{\uline{cf 5. v. p \uline{5bis}}}

\deleted{\unsure{5}}6. Nous allons étudier de façon plus détaillée un cas où les conditions (4.4.) sont satisfaites. Pour ceci, partons encore du cas particulier $M^*$, \inserted{$E^*$ sous-algèbre graduée de $\operatorname{End}(M^*)$,} et supposons qu'on ait un
\[
 L\in E^{(2)}
 \tag*{(6.0)}
\]
qui ait la propriété que
% archive sheet number: 6
\scanpage{8}\mspage{5 bis}

5. On peut aussi donner une condition d'\uline{existence} des $\pi_i$, liée aux conditions (4.4). Pour ceci, on suppose que $\unsure{E|\mathcal E}$ est contenu dans un sur-anneau gradué $F$ [p. ex. \inserted{l'anneau} des endom. d'un $M^*$ comme dans 1.] pour lequel les $\pi_i$ satisfont (1.1.) et les relations correspondantes à (1.3) (avec $E$ remplacé par $F$) sont satisfaites, \inserted{et (3.0) : $\pi_i=0$ si $i\notin[0,2n]$} Supposons que $\forall i$, la condition $(4.4)_i$ soit satisfaite, (avec $v_i,w_i\in E$ bien sur). On va en conclure que
\[
 \pi_i\in E\qquad\text{pour tt }i
 \tag*{(5.1)}
\]
\deleted{\unsure{Raisonnons prouvons en fait|Raisonnons Supposons provisoirement}} \inserted{supposons prouvé \deleted{\unsure{que}} déjà} $\pi_\alpha\in E$ si $\alpha\notin[i,2n-i]$, où $i<n$, et prouvons que l'on peut remplacer $i$ par $i+1$, \inserted{i.e $\pi_i$ et $\pi_{2n-i}\in E$.} Trivial si $i=n-1$, supposons donc $i\le n-2$.

Noter que, avec les notations du 3., \deleted{\unsure{on peut}} \inserted{on aura $\varphi_i,\psi_{2n-i}\in E$, \deleted{\unsure{d}} et} \deleted{\unsure{remplacer} $v_i$, \inserted{$w_i$} par} \inserted{donc} \deleted{\illegible{1}} $\psi_{2n-i}v_i\varphi_i,\ \varphi_iw_i\psi_{2n-i}\in E$, i.e
\[
 \left\{\begin{aligned}
 \inserted{v' =}\,\psi_{2n-i}v_i\varphi_i&=\pi_{2n-i}v_i\pi_i\in E\\
 w'_i=\varphi_i w_i\psi_{2n-i}&=\pi_i w_i\pi_{2n-i}\in E
 \end{aligned}\right.
\]
\deleted{et} \uline{co}mme \deleted{\unsure{on a}}
\[
 w'_iv'_i=\pi_i w_i\pi_{2n-i}v_i\pi_i
 \overset{\inserted{\text{car }w_i\pi_{2n-i}=\pi_i w_i}}{=}
 \pi_i\pi_i\unsure{W_i|w_i}v_i\pi_i
 \overset{\inserted{\text{car }w_iv_i\pi_i=\pi_i}}{=}\pi_i^3=\pi_i .
\]
et de $=$ \deleted{\unsure{$=\pi_i$}} $v'_i w'_i=\pi_{2n-i}$, on conclut que $\pi_i\in E$, $\pi_{2n-i}\in E$.
% archive sheet number: 7
\scanpage{9}\mspage{6}

\[
 L^{n-i}\mid M^i : M^i\longrightarrow M^{2n-i}\quad\text{est un isom.}
 \tag*{$(6.1)_i$}
\]
Ceci implique déjà $(4.1)_i$ et $(4.2)_i$ donc $(3.1)_i$ et $(3.2.)_i$, \inserted{donc l'\uline{unicité} des $\pi_i$ en termes de la structure graduée de $\unsure{E|\mathcal E}^*$.} Supposons de plus

$(6.2)_i$ L'inverse de l'isom (6.1) \deleted{\unsure{dans $E$|$E$}\illegible{1}} \inserted{est} induit par un \unsure{élmt|élt} de $E^{-(2n-2i)}$.

\deleted{Ceci} L'ensemble des deux conditions (6.1) et (6.2). équivaut [en termes intrinsèques à $E^*$, muni de \deleted{\unsure{$E$}} $L\in E^2$, si on suppose déjà les $\pi_i\in E^*$]
\[
 \left\{\begin{aligned}
 &\exists\,\unsure{W_i|w_i}\in E^{-(2n-2i)}\quad\text{tel qu'on ait}\\
 &(w_i L^{n-i}-1)\pi_i=0,\qquad (L^{n-i}w_i-1)\pi_{2n-i}=0
 \end{aligned}\right.
 \tag*{$(6.3.)_i$}
\]
Cette condition à son tour (indépendamment de la donnée préliminaire d'un $M^*$) implique $(4.4)_i$, \deleted{et on en déduit $(3.1)_i$ et $(3.2)_i$} donc \inserted{en vertu des n\textsuperscript{os} 4, 5} \inserted{$\pi_i\in E$ et} l'\uline{unicité} des $\pi_i$ en termes des conditions \inserted{(1.1) et} (1.3).

Revenons à $M^*$, et $L$ satisfaisant à $(6.1)_i$. Nous avons vu comment \deleted{\unsure{la}} cette condition implique la \uline{décomposition de} $M^*$ en termes \deleted{\unsure{nou}} de morceaux primitifs
\[
 \left\{\begin{aligned}
 M^i&\simeq P^i+LP^{i-2}+L^2P^{i-4}\ \unsure{-|\cdots}\\
 M^{2n-i}&\simeq L^{n-i}P^i+L^{n-i+1}P^{i-2}\ \unsure{-|\cdots}
 \end{aligned}\ \right|\quad\text{si }i\le n
 \tag*{(6.4)}
\]
On désigne par $\pi_{i,\alpha}$ la projection de $M^i$ sur $L^\alpha P^{i-2\alpha}$, par $\pi_{2n-i,\alpha}$ celle de $M^{2n-i}$ sur $L^{\unsure{ni+\alpha|n-i+\alpha}}P^{i-2\alpha}$.
% archive sheet number: 8
\scanpage{10}\mspage{7}

On a aussi l'hom $\Lambda$ de degré $-2$, défini par
\[
 \left\{\ \cdots\cdots\ \right. \tag*{(6.5)}
\]

\uline{Théorème} Les conditions suivantes \inserted{(i) à (iii)} sont équivalentes

(i) Les conditions \inserted{$(6.1)_i$ et} $(6.2)_i$ pour tt $i$, et les $\pi_i\in \unsure{E|\mathcal E}$.

\deleted{[\unsure{La graduation} de $E^*$ peut être définie par des $\pi_i$, et pour un \deleted{\unsure{quel}} \unsure{tel choix}]} \deleted{(\unsure{ii}) \unsure{la condition} $(6.3)_i$ pour tt $i$ [avec $\pi_i$ remplacés par $\pi_i$ \unsure{\ldots}]} \deleted{(\unsure{ii})}

\inserted{(ii) \deleted{\unsure{les} \illegible{1} et} On a la condition $(6.3)_i$ pour tout $i$}

(iii) \deleted{\unsure{L'hom}} \inserted{On a} \inserted{$(6.1)_i$ pour tt $i$, (d'où $\Lambda$) et}
\[
 \underline{\Lambda\in E^{(-2)}}\tag*{(6.6)}
\]

\unsure{(iv)} Ces conditions impliquent de plus

(iv) Les $\pi_{i,\alpha}$ sont dans $E$, \deleted{les $\pi_i$ sont dans $E$.}

(v) Les $\pi_i\in E$, \inserted{et ils} sont définis \inserted{(en termes intrinsèques à $E$ gradué)} comme formant \deleted{les seuls projecteurs} la seule famille d'élémts de $E$ satisfaisant (1.1.) et (1.3) (a) ou b), au choix \ldots, \deleted{\unsure{E}} ou encore (1.4)\ldots).

\deleted{(vi) L'élèmt $\Lambda$ de $E^{(-2)}$ est uniquement déterminé, \deleted{\unsure{par les}} (en termes de \inserted{l'anneau gradué} $E^*$ et l'élèmt $L\in E^{(\unsure{2|1})}$, par les conditions suivantes :
\[
 \left\{\begin{aligned}
 L\Lambda&=1-\sum_{0\le i\le n}\pi_{i,0}\\
 \Lambda L&=1-\sum_{n\le i\le 2n}\pi_{i,0}
 \end{aligned}\right.
 \tag*{(6.7)}
\]}

\uline{Dém} L'équivalence de (i) et (ii) est claire, et le fait que c'est impliqué par (iii) (en prenant $w_i=\Lambda^{\deleted{\unsure{(}}n-i\deleted{\unsure{)}}}$). \deleted{\unsure{Il}} \inserted{Reste à} prouver que (i), (ii) implique (iii) à (\unsure{vi|v}). On \deleted{\unsure{sait déjà}} \inserted{a déjà signalé} que ça implique \deleted{\unsure{(6.3)}} (v).
% archive sheet number: 9
\scanpage{11}\mspage{8}

Prouvons (iv). Supposons prouvé
\[
 \pi_{j\alpha}\in \unsure{E|\mathcal E}\qquad\text{pour }j\notin[i,2n-i]
 \tag*{(6.7)}
\]
où $i\le n$, et prouvons que
\[
 \deleted{\pi_{j\alpha}\in E\quad\text{pour } j\in}
\]
\[
 \pi_{j\alpha}\in E\quad\deleted{\text{pour}}\ \text{si }j=i\ \text{ou }j=2n-i,
 \tag*{$(6.7')$}
\]
ce qui prouvera (iv) \inserted{par récurrence ascendante sur $i\le n$.} Pour prouver 6.8. on procède par récurrence \deleted{\unsure{sur}} descendante sur $\alpha$, en supposant qu'on a
\[
 \deleted{\unsure{\pi_{j\alpha}}}\quad \pi_{i,\alpha'}\in E,\quad\pi_{2n-i,\alpha'}\in E\quad\text{si }\alpha'>\alpha,
 \tag*{(6.8)}
\]
et on en \unsure{conclura|conclure} que
\[
 \pi_{i,\alpha}\in E,\quad\pi_{2n-i,\alpha}\in E .
 \tag*{$(6.8')$}
\]
\deleted{Cela \deleted{\unsure{étant}} prouvera bien (6.7'), puisque on peut démarrer la récurrence ($\pi_{i,\alpha}=\pi_{2n-i,\alpha}=0$ si $\alpha$ grand). Si $\alpha=0$, on gagne grâce à la relation
\[
 \sum_\alpha\pi_{i,\alpha}=\pi_i\quad\text{donc}\quad\pi_{i0}=\pi_i-\sum_{\alpha\ge1}\pi_{i,\alpha},
 \tag*{(6.9)}
\]
donc il suffit de regarder le cas $\alpha\ge1$.}

Posons
\[
 q_{i,\alpha}=1-\sum_{\alpha'>\alpha}\pi_{i,\alpha'}=\sum_{\alpha'\le\alpha}\pi_{i,\alpha}
\]
donc grâce à (6.8) on a
\[
 q_{i,\alpha}\in E\ \deleted{\unsure{E}}
\]
% archive sheet number: 10
\scanpage{12}\mspage{9}

Alors on voit que
\[
 L^{n-i+\alpha}q_{i\alpha}: M^i\unsure{-|\to}M^{2n-i+2\alpha}
\]
est nul sur les $\deleted{\unsure{P}}L^{\alpha'}P^{i-2\alpha'}$ pour $\alpha'\ne\alpha$, et que sur $L^\alpha P^{i-2\alpha}$, \deleted{\unsure{c'est}} son composé avec $P^{i-2\alpha}\xrightarrow{L^\alpha} L^\alpha P^{i-2\alpha}$ n'est autre que l'isom
\[
 L^{n-i-2\alpha}:P^{i-2\alpha}\longrightarrow L^{n-i+2\alpha}P^{i-2\alpha}\subset M^{2n-i+2\alpha}.
\]
Il s'ensuit aussitôt qu'on a
\[
 \pi_{i,\alpha}=L^\alpha w_{i-2\alpha}L^{n-i+\alpha}q_{i\alpha}
 =L^\alpha w_{i-2\alpha}L^{n-i+\alpha}(1-\sum_{\alpha'>\alpha}\pi_{i,\alpha'}\unsure{)}
 \tag*{(6.9)}
\]
ce qui donne bien $\pi_{i\alpha}\in\unsure{E|\mathcal E}$. \deleted{Cela prouve (iv).} Si $i<n$, il faut \unsure{égalt} $\pi_{2n-i,\alpha}\in E$. On

\deleted{Reste à prouver \inserted{(iii)} $\Lambda\in E$. Il suffit de \deleted{\unsure{prouv.}} \unsure{poser}
\[
 q_{2n-i,\alpha}=1-\sum_{\alpha'>\alpha}\pi_{2n-i,\alpha'}
\]
et on voit alors que
\[
 \deleted{\unsure{q_{2n-i,\alpha}|\pi_{2n-i,\alpha}}} L^\alpha q_{2n-i,\alpha}:M^{2n-i}\unsure{-|\to}M^{2n-i+2\alpha}
\]}
le déduit \unsure{p\textsuperscript{r} ex.|par ex.} de la formule
\[
 \pi_{2n-i,\alpha}=L^{n-i}\pi_{i,\alpha}w_i\deleted{\unsure{\pi_{2n-i}}}
 \tag*{(6.10)}
\]
Cela prouve \inserted{(6.8') donc} (iv). Reste à prouver (iii) $\Lambda\in E$.

Comme on a
\[
 \Lambda=\sum_{\substack{j,\alpha\\\deleted{\unsure{0\le j\le 2n}}}}\deleted{\unsure{\{}}\Lambda\pi_{j,\alpha}
 \deleted{\unsure{+\sum_{\substack{\alpha>0\\0\le i\le n}}\Lambda\pi_{i,\alpha}\}}}
 \tag*{(6.11)}
\]
il suffit de prouver que $\Lambda\pi_{i\alpha}\deleted{\unsure{\Lambda\pi_{2n-i,\alpha}}}\in E$.
% archive sheet number: 11
\scanpage{13}\mspage{10}

Or si $i\le n$, on aura
\[
 \left\{\ \Lambda\pi_{i0}=0,\quad \Lambda\pi_{i\alpha}=L^{\alpha-1}\unsure{W_{i-2\alpha}|w_{i-2\alpha}} L^{n-i+\alpha}\pi_{i\alpha}\quad\text{si }\alpha\ge1,\right.
 \tag*{(6.12)}
\]
\deleted{et si $j>n$, on aura}
\[
 \deleted{\Lambda\pi_{j,\alpha}=L}\tag*{\deleted{(6.13)}}
\]
\[
 \Lambda\pi_{2n-i,\alpha}=L^{n-i+\alpha-1}w_{i-2\alpha}\deleted{\unsure{\pi_{2n-i,\alpha}}}L^\alpha\pi_{2n-i,\alpha}\quad(\text{si }i\le n-1)
 \tag*{(6.13)}
\]
Les conditions envisagées impliquent évidemment (iv), : \unsure{prouvons|prouver} que (iv) implique (ii). Il suffit de prouver en fait que l'on a alors $\pi_i=\pi'_i$. \uline{Démonstration à fournir}.

\uline{Remarque} Voici un système de conditions déterminant de façon unique les éléments $\Lambda\in\unsure{E|\mathcal E}^{(2)}$, $\pi_{i,\alpha}\in E^{(0)}$ : en termes de $L$, et des $\pi_i^{\deleted{\unsure{*}}}$ (dont nous avons déjà une caractérisation intrinsèque par ailleurs, en termes de l'anneau gradué $E$) :

\margin{\deleted{\unsure{$\forall i$}}}
\[
 \left\{\begin{aligned}
 \text{\textcircled{\scriptsize 1}}\quad&(\Lambda^{n-i}L^{n-i}-1)\pi_i=0\qquad
 \text{\textcircled{\scriptsize 2}}\ \deleted{\unsure{\pi_{2n-i}}}(L^{n-i}\Lambda^{n-i}-1)\deleted{\unsure{\pi_i}}\pi_{2n-i}=0\\
 \text{\textcircled{\scriptsize 3}}\quad&\pi_{i,\alpha}=L^\alpha\deleted{\unsure{w_{i-2\alpha}}}\Lambda^{n-i+2\alpha}L^{n-i+\alpha}(1-\sum_{\alpha'>\alpha}\pi_{i,\alpha'})\\
 \text{\textcircled{\scriptsize 4}}\quad&\pi_{2n-i,\alpha}=L^{n-i}\pi_{i,\alpha}\Lambda^{n-i}\\
 &\deleted{\Lambda=\sum_{j,\alpha}\Lambda\pi_{j,\alpha}}\\
 \text{\textcircled{\scriptsize 5}}\quad&\Lambda\pi_{i,0}=0\ \unsure{\alpha}\\
 \text{\textcircled{\scriptsize 6}}\quad&\Lambda\pi_{i\alpha}=L^{\alpha-1}\Lambda^{n-i+2\alpha}L^{n-i+\alpha}\pi_{i\alpha}\quad\text{si }\alpha\ge1\\
 \text{\textcircled{\scriptsize 7}}\quad&\Lambda\pi_{2n-i,\alpha}=L^{n-i+\alpha-1}\Lambda^{n-i+2\alpha}L^\alpha\pi_{2n-i,\alpha}\quad\text{si }\unsure{\alpha|i}\le n-1\\
 \text{\textcircled{\scriptsize 8}}\quad&\Lambda=\sum_{j,\alpha}\Lambda\pi_{j,\alpha}
 \end{aligned}\right.
 \tag*{(6.14)}
\]

(dans ttes ces formules, on suppose $i\le n$).

Dans ces huit formules, on peut d'ailleurs remplacer les quatre dernières par les suivantes \inserted{plus simples} :
% archive sheet number: 12
\scanpage{14}\mspage{11}

\[
 \left\{\begin{aligned}
 L\Lambda&=1-\sum_{0\le i\le n}\pi_{i,0}\\
 \Lambda L&=1-\sum_{n\le j\le 2n}\pi_{\deleted{\unsure{2i}}j,0}
 \end{aligned}\right.
 \tag*{(6.15)}
\]
\deleted{Quant aux $\pi_{j,0}$, ils sont déterminés par les}

\uline{Question} Donner un système de \deleted{\unsure{générateurs}} relations en les $\Lambda,L$ exclusivement, \deleted{\unsure{puis}} des relations exprimant les $\pi_i$ comme polynômes non comm en $L,\Lambda$ ou qui définissent $\Lambda$ en termes de $L$ et des $\pi_i$ \ldots
% slanted
\margin{cf n\textsuperscript{o} 7 pour l'anneau $\Phi$ : système de relations \ldots}

7. Revenant à un anneau gradué abstrait $\unsure{E|\mathcal E}$, muni d'un
\[
 L\in E^{(2)},
\]
on veut abstraire la situation du n\textsuperscript{o} précédent, sans utiliser $M^*$, en termes des propriétés du seul anneau \deleted{\unsure{$E$}} gradué $E^*$ muni de $L\in E^{(2)}$.

Partons pour ceci d\deleted{\unsure{e}}'un module libre $\mathbb{M}_{\unsure{0|\circ}}$ de rang $2n$ sur $\mathbb Z$, muni d'une 2-forme alternée non dégénérée \inserted{(sur $\mathbb Z$ !)} $\unsure{\xi|\zeta}_0\in\bigwedge^2\mathbb M_0$, \inserted{$\xi^*\in\bigwedge^2\mathbb M_0^\vee$ la forme duale,} soit \deleted{\unsure{$\xi$}} $\mathbb M_0^*=\bigwedge\mathbb M_0$, et considérons sur $M_0^\alpha$ les opérateurs
\[
 \mathbb L x=\xi_0x,\qquad \mathbb{\Lambda}_0x=\xi_0^*\mathbin{\lrcorner}x .
\]
et considérons l'algèbre
\[
 \Phi_0=\mathbb Z[\mathbb L_0,\mathbb{\Lambda}_0]\subset\operatorname{End}_{\mathbb Z}(M_0^*)\qquad\text{anneau \unsure{engendré|gradué}.}
\]
Noter que $\operatorname{End}(\mathbb M_0^*)$, \inserted{opérant sur $\mathbb M_0^*$,} muni de $\mathbb L$, satisfait à la condition de (6.1) \unsure{(N)|($\forall$)|($\omega$)}, \deleted{\unsure{donc}} et \deleted{\unsure{$\xi$}} $\mathbb{\Lambda}_0$ est l'opération associée définie par (6.5).
% archive sheet number: 13
\scanpage{15}\mspage{12}

\deleted{\unsure{Noter que}} $\Phi_{\unsure{0|\circ}}$ est un module libre sur $\mathbb Z$, engendré par les
\[
 \varphi_i^{\beta,\alpha}\deleted{\unsure{\pi_{i,\alpha}}}
 \qquad\begin{aligned}&0\le i\le n\\&0\le\alpha,\beta\le2n-i\end{aligned}
\]
\[
 \deleted{\left\{ L^j\pi_{i,\alpha}\quad\varphi_{i,\alpha}^{j,\beta}\ \unsure{\pi_{i,\alpha}}\quad\text{pour }
 \left\{\begin{aligned}
 &\unsure{\text{avec }}0\le i\le 2n,\ 0\le j\le2n,\ \alpha,\beta\ge0\\
 &i-2\alpha=j-2\beta\ge0\\
 &\unsure{0\le i<2}
 \end{aligned}\right.\right.}
\]
% slanted
\margin{$\mathbb Z$ peut être remplacé par n'importe quel anneau}
\deleted{$\varphi_{i,\alpha}^{j,\beta}$ est nul sur les composantes} définis par
\[
 \left|\begin{aligned}
 \varphi_i^{\beta,\alpha}&=L_0^{\beta-\alpha}\pi_{i+2\alpha,\alpha}\quad\text{si }\alpha\le\beta\\
 \varphi_i^{\beta,\alpha}&=\Lambda_0^{\alpha-\beta}\pi_{i+2\alpha,\alpha}\quad\text{si }\alpha\ge\beta
 \end{aligned}\right.
\]
\deleted{\[
 \left\{\begin{aligned}
 \varphi_{i,\alpha}^{j,\beta}&=L_0^{\beta-\alpha}\pi_{i,\alpha}=\pi_{j,\beta}L_0^{\beta-\alpha}\quad\text{si }\alpha\le\beta\\
 \varphi_{i,\alpha}^{j,\beta}&=\Lambda_0^{\alpha-\beta}\pi_{i,\alpha}=\pi_{j,\beta}\Lambda_0^{\alpha-\beta}\quad\text{si }\alpha\ge\beta
 \end{aligned}\right.
\]}
i.e. on a une base formée des
\[
 \varphi_{i'}^{\beta',\alpha'}\varphi_i^{\beta,\alpha}
 \quad\deleted{\unsure{\varphi_{i,\alpha}^{j,\beta}\varphi_{i',\alpha'}^{j',\beta'}}}
 =\left|\begin{aligned}
 &0\ \text{si }\deleted{\unsure{j\ne i'\text{ ou }2\alpha\ne\beta}}\inserted{i\ne i'\text{ ou }\beta\ne\alpha'}\\
 &\varphi_{\unsure{i|i'}}^{\beta',\alpha}\quad \text{sinon}
 \end{aligned}\right.
\]
\deleted{\[
 \left\{ L^\beta\pi_{i,\alpha}\qquad
 \left|\begin{aligned}
 &0\le i\le\unsure{2n|n},\quad\alpha\ge0\quad\text{tel que }i-2\alpha\ge0\\
 &0\le2\alpha\le i\\
 &0\le\beta
 \end{aligned}\right.\right.
\]}
% slanted
\margin{Les $\varphi_i^{\beta\alpha}$ et leurs multiplications devraient être donnés au n\textsuperscript{o} 6}

La multiplication des élémts de cette base ne dépend pas du choix de $\mathbb M_0,\unsure{\xi|\zeta}_0$ (ni du fait qu'on est sur l'anneau $\mathbb Z$ \ldots), donc l'anneau $\Phi_0$ est bien déterminé en termes de l'entier \inserted{pair} $2n$.

\uline{Proposition} Revenant à la donnée de $M^*$ comme au n\textsuperscript{o} 1, et de $\unsure{E|\mathcal E}$ ss-\deleted{algèbre}\inserted{anneau} graduée de \deleted{\unsure{$E^*$}} $\operatorname{End}(M^*)$. Soit $L\in E^{(\unsure{2|1})}$. Conditions équivalentes

(i) L'endom $L$ de $M^*$ satisfait la condition de \deleted{\unsure{Hodge}} Lefschetz, et l'opérateur $\Lambda$ associé est dans $E$. (i.e on est \unsure{sous|dans} les conditions du th. du n\textsuperscript{o} 6).

(ii) Il existe un hom d'anneaux gradués
\[
 \varphi:\Phi_0\longrightarrow E
\]
% archive sheet number: 14
\scanpage{16}\mspage{13}

appliquant $\mathbb L_{\unsure{0|\circ}}$ en $L$, \uline{et les} $\mathbb{\pi}_i$ \uline{en les} $\pi_i$ \deleted{\unsure{définis}} \deleted{\unsure{satisf} la graduation de $\unsure{E|\mathcal E}$ par (1.3).} (donc $\pi_i\in E$)

\margin{condition essentielle !}
% slanted
\margin{\deleted{(iii.) Il existe \deleted{\unsure{une représentation}} un \unsure{$E$-module} \uline{fidèle} $M^*$ de $E^*$, \deleted{\unsure{fiel}} que les conditions du th\ldots\ du n\textsuperscript{o} 6 soient satisfaites (pour la graduation envisagée) (iv) pour \uline{tt} module fidèle \ldots}}

De plus, sous les conditions, l'hom annoncé dans (ii) est unique, \deleted{appliquant} et avec les notations du n\textsuperscript{o} 6, il transforme $\mathbb{\Lambda}_0$ en $\Lambda$, les \deleted{\unsure{$\pi_\alpha$}} $\mathbb{\pi}_{i\alpha}$ en $\pi_{i\alpha}$, $\mathbb{\pi}_i$ en $\pi_i$, \deleted{\illegible{2}} \inserted{$\varphi_i^{\beta,\alpha}$ en $\varphi_i^{\beta,\alpha}$}.

\uline{Dém} (i) $\Rightarrow$ (ii) Résulte de la table de multiplication des $\varphi_i^{\beta,\alpha}$ donnée au n\textsuperscript{o} 6, et dudit th. donné au n\textsuperscript{o} 6.

(ii) $\Rightarrow$ (i) \deleted{Soient $\pi'_i$ les images des $\mathbb{\pi}_i$. On aura $L^\alpha\pi'_i=\pi'_{i+\alpha}L^\alpha$ \unsure{voit|fait} alors que les $\pi'_i$ définissent la graduation de $E$, puisque ils sont \unsure{de|du} degré 0. Je dis que $\pi'_i=\pi_i$ pour tt $i$. Cf Proposition ci-dessous pour ce point, appliquée à $\Phi\unsure{-|\to}\operatorname{End}(M^*)$.}

Comme on a, pour tt $i\le n$
\[
 (\Lambda_0^{n-i}\mathbb L_0^{n-i}-1)\pi_i=0,\quad
 (\mathbb L_0^{n-i}\Lambda_0^{n-i}-1)\pi_{2n-i}=0
\]
on en conclut, appliquant $\varphi$ (et posant $\Lambda'=\varphi(\mathbb{\Lambda})$)
\[
 ({\Lambda'}^{n-i}L^{n-i}-1)\pi_i=0,\quad (L^{n-i}{\Lambda'}^{n-i}-1)\pi_{2n-i}=0
\]
Donc la condition (ii) du th. du n\textsuperscript{o} 6 est vérifiée, donc aussi la condition (i).

Pour l'unicité de $\varphi$, \deleted{\unsure{et donc sa proposition en}} \deleted{\unsure{et les autres}} il suffit de voir que $\varphi(\mathbb{\Lambda})=\Lambda$, et \inserted{pour avoir $\varphi(\mathbb{\pi}_{i,\alpha})=\pi_{i,\alpha}$,} \deleted{\unsure{au n\textsuperscript{o} 2 en réalité,}} compte tenu de $\varphi(\mathbb{\pi}_i)=\pi_i$, \deleted{\unsure{puisque des} \illegible{1}. Pour $\varphi(\mathbb{\Lambda}$} \deleted{\unsure{du fait}} \inserted{on utilise} la caractérisation \inserted{simultanée} de $\Lambda$ \inserted{et des $\pi_{i,\alpha}$} par les formules (6.\unsure{14|15}).
% archive sheet number: 15
